\documentclass[11pt,a4paper]{article}
\usepackage{isabelle,isabellesym}
\usepackage{fontspec}
\usepackage{polyglossia}
\usepackage{amssymb}
\usepackage{amsmath}

\setdefaultlanguage{russian}
\setotherlanguage{english}

% Нарратив (root + isamarkuptext): Charis SIL.
% Леммы / определения Isabelle (\isabellestyle{it}): Minion Pro.
% Mono: DejaVu --- символы и код.
\setmainfont{Charis SIL}[Ligatures=TeX]
\setsansfont{DejaVu Sans}[Ligatures=TeX]
\setmonofont{DejaVu Sans Mono}[Scale=0.92]
\newfontfamily\cyrillicfont{Charis SIL}[Ligatures=TeX]
\newfontfamily\cyrillicfontsf{DejaVu Sans}[Ligatures=TeX]
\newfontfamily\cyrillicfonttt{DejaVu Sans Mono}[Scale=0.92]
\newfontfamily\isabellefont{Minion Pro}[Ligatures=TeX]

% this should be the last package used
\usepackage{pdfsetup}

\urlstyle{rm}
\isabellestyle{it}
% Переопределяем стили Isabelle на Minion Pro (нарратив остаётся Charis/\rmfamily).
\def\isastyle{\small\isabellefont\itshape}
\def\isastyleminor{\isabellefont\itshape}
\def\isastylescript{\footnotesize\isabellefont\itshape}

\begin{document}

\title{Учебная модель передачи capability в стиле Zircon\\
в Isabelle/HOL}
\author{redbeardster@gmail.com}
\maketitle

\begin{abstract}
В работе представлена \emph{учебная} формальная модель передачи прав
доступа (capability) в духе микроядра Zircon (основа Fuchsia OS) на языке
Isabelle/HOL. Мы описываем процессы с таблицами дескрипторов (handles),
каналы с очередями capabilities, сокращение прав при
\texttt{zx\_channel\_write} и проверку права записи при \texttt{vmo\_write}.
Доказаны: невозможность эскалации прав при передаче, сохранение ссылки на
объект, необходимость \texttt{ZX\_RIGHT\_TRANSFER}, а также целостность
данных VMO (запись только при наличии \texttt{ZX\_RIGHT\_WRITE}).

Модель намеренно упрощена и \emph{не} является верификацией реального ядра
Zircon и \emph{не} даёт оснований утверждать корректность TCB Fuchsia.
\end{abstract}

\section{Введение}

Zircon использует capability-модель: пользовательское пространство обращается
к объектам ядра через \emph{handles}~--- локальные идентификаторы, к каждому
из которых привязана маска прав. Handles передаются между процессами через
каналы; при передаче права могут только уменьшаться.

Цель этой работы~--- зафиксировать небольшой фрагмент этой идеи в Isabelle/HOL
так, чтобы заявленные свойства \emph{реально следовали} из определений и
проверялись машиной. Мы явно отделяем:
\begin{itemize}
  \item \textbf{постулаты модели}~--- определения и упрощения, которые мы
        принимаем как аксиоматику учебной системы;
  \item \textbf{доказанные утверждения}~--- леммы и теоремы, автоматически
        проверенные Isabelle.
\end{itemize}

\section{Что мы постулируем}

Ниже перечислены решения модели, которые \emph{не доказываются}, а задаются
определениями. От них зависят все дальнейшие теоремы.

\subsection{Онтология}

\begin{enumerate}
  \item Идентификаторы (\texttt{handle\_id}, \texttt{obj\_id},
        \texttt{process\_id})~--- натуральные числа, а не
        32-битные слова Zircon.
  \item Объект ядра~--- пара $\langle\texttt{тип}, \texttt{id}\rangle$; глубокая
        семантика объектов (кроме абстрактного содержимого VMO) не моделируется.
  \item Handle~--- ссылка на объект плюс множество прав.
  \item Глобальное состояние~--- таблицы handles процессов, счётчики id,
        очереди каналов (ключ~--- \texttt{obj\_id} объекта
        \texttt{Channel}) и склад содержимого VMO.
  \item Канал~--- kernel object: доступ к очереди только через handle на
        endpoint с \texttt{WRITE}/\texttt{READ}.
\end{enumerate}

\subsection{Правила операций}

\begin{enumerate}
  \item Сокращение прав~--- пересечение множеств:
        $R_{\mathrm{actual}} = R_{\mathrm{req}} \cap R_{\mathrm{src}}$.
  \item \texttt{zx\_channel\_write} требует handle на канал с
        \texttt{ZX\_RIGHT\_WRITE} и payload с \texttt{ZX\_RIGHT\_TRANSFER};
        кладёт capability в очередь и \emph{удаляет} payload-handle
        отправителя (move).
  \item \texttt{transfer\_via\_channel} работает только на \emph{пустом}
        канале, чтобы получатель взял именно только что записанную capability.
  \item \texttt{vmo\_write} проверяет маску при каждом вызове: без
        \texttt{WRITE} возвращается \texttt{None}, состояние не меняется.
  \item При \texttt{create\_vmo} можно указать любые начальные права
        (политика minting вне scope).
\end{enumerate}

\subsection{Чего мы сознательно не моделируем}

Параллелизм, полный syscall API, refinement к C++, полный набор rights
Zircon, политика minting при создании объектов.

\section{Операции модели (кратко)}

\begin{center}
\begin{tabular}{ll}
\texttt{create\_channel\_pair} & создать Channel и handles двум процессам \\
\texttt{create\_vmo} & создать VMO и handle в таблице процесса \\
\texttt{zx\_channel\_write} & move capability в канал через endpoint-handle \\
\texttt{zx\_channel\_read} & забрать capability в таблицу получателя \\
\texttt{transfer\_via\_channel} & write+read на пустом канале \\
\texttt{vmo\_write} & изменить содержимое VMO при наличии WRITE
\end{tabular}
\end{center}

Проверка прав происходит и при передаче (пересечение + TRANSFER), и при
мутирующем syscall (\texttt{vmo\_write}).

\section{Доказанные свойства: путеводитель по леммам}

\subsection{Алгебра прав}

\begin{description}
  \item[\texttt{reduce\_rights\_subset}]
        $R_{\mathrm{new}} \cap R_{\mathrm{old}} \subseteq R_{\mathrm{old}}$.
        Базовый факт: пересечение не добавляет прав.
  \item[\texttt{reduce\_rights\_idempotent}]
        Два сокращения подряд = одно с пересечением запросов.
\end{description}

\subsection{Характеризация write/read}

\begin{description}
  \item[\texttt{zx\_channel\_write\_success}]
        Успешный write $\Rightarrow$ были WRITE на канал и TRANSFER на
        payload; в очередь добавлена capability (в т.\,ч.\ с пустой маской);
        payload-handle отправителя удалён (move).
  \item[\texttt{zx\_channel\_read\_success}]
        Успешный read $\Rightarrow$ capability с головы очереди установлена
        под свежим handle id у получателя.
  \item[\texttt{lookup\_after\_set\_table} / \texttt{lookup\_after\_alloc\_hid}]
        Технические факты о корректности lookup после обновлений.
\end{description}

\subsection{Шлюз TRANSFER}

\begin{description}
  \item[\texttt{channel\_write\_requires\_transfer}]
        Успех write $\Rightarrow$ у источника есть TRANSFER.
  \item[\texttt{channel\_write\_fails\_without\_transfer}]
        Нет TRANSFER $\Rightarrow$ write = \texttt{None}.
\end{description}

\subsection{Передача без эскалации}

\begin{description}
  \item[\texttt{transfer\_via\_channel\_cap}]
        Ключевая лемма: получатель имеет ровно
        $\langle\mathrm{obj}_{\mathrm{src}},\;
        R_{\mathrm{req}}\cap R_{\mathrm{src}}\rangle$.
  \item[\texttt{transfer\_via\_channel\_no\_escalation}]
        Права получателя $\subseteq$ права отправителя.
  \item[\texttt{transfer\_via\_channel\_preserves\_object}]
        Тот же \texttt{kernel\_obj}.
  \item[\texttt{channel\_transfer\_no\_escalation}]
        Теорема: не существует права, появившегося «из ниоткуда» у
        получателя.
\end{description}

\subsection{Целостность VMO и общий запрет записи}

\begin{description}
  \item[\texttt{vmo\_write\_denied\_without\_write\_right}]
        Специально для VMO: валидный VMO-handle без WRITE $\Rightarrow$
        запись запрещена. Это узкий инвариант целостности VMO.
  \item[\texttt{vmo\_write\_denied\_means\_no\_new\_state}]
        Шире: для \emph{любого} валидного handle без WRITE вызов
        \texttt{vmo\_write} даёт \texttt{None} --- в том числе если это
        handle на Process/Thread и т.п., у которого WRITE нет в маске
        (или объект вовсе не VMO). Не ошибка, а более общий факт о шлюзе
        WRITE у операции записи.
  \item[\texttt{vmo\_write\_succeeds\_iff\_write\_right}]
        Для VMO-handle успех $\Leftrightarrow$ есть WRITE и точечное
        обновление store.
  \item[\texttt{vmo\_write\_updates\_only\_with\_write\_right}]
        Успех $\Rightarrow$ был WRITE и store обновлён ожидаемо.
\end{description}

\subsection{Прочее}

\begin{description}
  \item[\texttt{channel\_write\_moves\_sender\_handle}]
        Write аннулирует payload-handle отправителя (move).
  \item[\texttt{example\_scenario\_true}]
        Полный сценарий A$\to$B с READ-only доказуемо истинен: B не может
        изменить содержимое VMO; у A payload-handle исчез.
  \item[\texttt{example\_empty\_rights\_ok\_true}]
        Передача с запросом $\emptyset$ допустима: получатель имеет пустую
        маску, объект тот же, \texttt{vmo\_write} отказывает.
\end{description}

\section{Ограничения (ещё раз)}

Модель полезна для обучения capability-системам и Isabelle/HOL. Она
\emph{не} доказывает безопасность производственного Zircon/Fuchsia. Любые
формулировки в духе «complete verification of the TCB» были бы неверны.

\clearpage

\tableofcontents

\clearpage

\parindent 0pt\parskip 0.5ex

% generated text of all theories
\input{session}

\clearpage

\section{Заключение}

Мы формализовали учебный фрагмент capability-модели:
\begin{enumerate}
  \item per-process handles с масками прав;
  \item передачу через канал с пересечением прав и проверкой TRANSFER;
  \item проверку WRITE на мутирующем syscall;
  \item машинно-проверенные доказательства перечисленных свойств.
\end{enumerate}

Постулаты модели заданы определениями; свойства безопасности выведены из них
и проверены Isabelle. Это шаг к пониманию capability-инвариантов, но не
замена верификации реального ядра.

\clearpage

\nocite{*}
\bibliographystyle{abbrv}
\bibliography{root}

\end{document}

%%% Local Variables:
%%% mode: latex
%%% TeX-master: t
%%% End:
